\documentclass[12pt]{article}
\usepackage[T1]{fontenc}
\usepackage[utf8]{inputenc}
\usepackage{lmodern}
\usepackage{textcomp}
\usepackage{enumitem}
\usepackage{geometry}
\geometry{margin=1in}
\begin{document}
\begin{center}
Answer Key - Sociology - Undergraduate Level Assessment 8 \
\small (Answers are ordered exactly as in the question paper.)
\end{center}

\section*{Multiple Choice}
\begin{enumerate}[label=\arabic*.]
\item \textbf{Answer:} B
\\ \textit{Options:} A) an increasing awareness of the world as a whole; B) the extended power of nation states; C) the destruction of distance through communications technologies; D) the stretching of social relations beyond national boundaries
\\ \textit{Note:} The extended power of nation states is not a feature of globalization because globalization typically involves the diminishing influence of individual nation states in favor of global networks and transnational entities.
\item \textbf{Answer:} A
\\ \textit{Options:} A) declining profits and rising unemployment; B) the eradication of the welfare state; C) rising divorce rates and the decline of the traditional family; D) an unfortunate twist in fashion sensibility
\\ \textit{Note:} The term 'crisis of the 1970 s' accurately describes the economic turmoil of that decade, characterized by declining profits in key industries and rising unemployment rates, which stemmed from factors such as oil shocks, inflation, and shifts in global economic dynamics.
\item \textbf{Answer:} C
\\ \textit{Options:} A) British-born second generation immigrants from the Asian subcontinent; B) White Americans who wanted to convert to Islam; C) African-Americans who felt excluded from the 'ethnic melting pot' in the USA; D) African-Caribbeans who lived in the inner cities and had a distinctive youth culture
\\ \textit{Note:} The Nation of Islam appealed to African-Americans who felt excluded from the 'ethnic melting pot' in the USA because it provided a sense of identity, community, and empowerment in response to systemic racism and marginalization experienced by Black individuals in American society.
\item \textbf{Answer:} A
\\ \textit{Options:} A) Functionalism; B) Conflict theory; C) Symbolic interactionism; D) Role theory
\\ \textit{Note:} Functionalism is the correct answer because it emphasizes how various components of society work together in a coordinated manner to maintain stability and social order, viewing society as an interdependent system.
\item \textbf{Answer:} C
\\ \textit{Options:} A) Labour Force Survey; B) General Household Survey; C) Fashion Sensibility Survey; D) Family Expenditure Survey
\\ \textit{Note:} The Fashion Sensibility Survey is not a regular national survey conducted by the British government, as it focuses on niche consumer interests rather than on broader social, economic, or demographic data typically gathered through official surveys.
\end{enumerate}

\section*{True / False}
\begin{enumerate}[label=\arabic*.]
\setcounter{enumi}{5}
\item \textbf{Answer:} True
\\ \textit{Options:} A) True; B) False
\\ \textit{Note:} The statement is true because Bell's (1973) concept of post-industrial society emphasizes the transition from manufacturing-based economies to those driven by services, knowledge, and information, reflecting a fundamental change in the nature of work and economic value.
\item \textbf{Answer:} True
\\ \textit{Options:} A) True; B) False
\\ \textit{Note:} Marx argued that capitalism leads to the commodification of labor, causing workers to become estranged from the products of their work, the act of production itself, and their fellow workers, ultimately undermining their sense of connection and fulfillment.
\item \textbf{Answer:} False
\\ \textit{Options:} A) True; B) False
\\ \textit{Note:} Bernstein argued that restricted codes, which are characterized by limited vocabulary and context-specific meanings, hindered children's ability to communicate effectively across diverse peer groups and contexts, as opposed to promoting effective communication.
\item \textbf{Answer:} True
\\ \textit{Options:} A) True; B) False
\\ \textit{Note:} The statement is true because a social stratum represents a distinct tier within the social hierarchy, characterized by individuals who share similar socioeconomic status, opportunities, and life experiences, which influences their access to resources and social mobility.
\item \textbf{Answer:} True
\\ \textit{Options:} A) True; B) False
\\ \textit{Note:} Frank's theory of dependency posits that 'underdeveloped' societies are trapped in a cycle of exploitation and economic dependency due to their historical and structural ties to wealthier nations, which extract resources and labor while perpetuating inequality and hindering local development.
\end{enumerate}

\section*{Long Form Response}
\begin{enumerate}[label=\arabic*.]
\setcounter{enumi}{10}
\item \textbf{Answer:} Berger and Luckmann's concept of the social construction of reality highlights how individuals actively negotiate and create shared definitions of their situations within society. This process involves interactions in which people communicate and reinforce common understandings and meanings, allowing them to navigate their social worlds effectively. As individuals adhere to these constructed realities, they often forget that these definitions are not inherent truths but rather fluid and subject to change. This dynamic illustrates the importance of social context in shaping perceptions and behaviors, emphasizing that reality is a product of collective agreement rather than an immutable external force. Ultimately, it showcases the active role individuals play in the ongoing construction of their social environments.
\\ \textit{Note:} correct because it emphasizes the active role of individuals in shaping and negotiating shared meanings through social interactions, illustrating that these constructed realities are not fixed truths but are continuously formed and reinforced within societal contexts.
\item \textbf{Answer:} The concept of new forms of community in sociology refers to the evolving social structures that emerge in response to technological advancements, globalization, and changes in social interaction patterns. These communities often thrive on shared interests, virtual connections, and inclusivity, fostering a sense of belonging among diverse members. However, sociological communities formed by unpopular lecturers do not fit this classification because they typically lack the essential elements of engagement, mutual support, and shared identity that characterize successful new communities. Instead, these groups often consist of individuals bound by obligation rather than genuine connection, leading to a disjointed and less cohesive experience. Consequently, such communities fail to cultivate the dynamic interactions necessary for a thriving sociological community.
\\ \textit{Note:} correct because it identifies that new forms of community are characterized by inclusivity and shared interests, which are absent in communities formed by unpopular lecturers, who often struggle to create a sense of belonging and engagement among members.
\end{enumerate}

\vfill
\noindent\textit{End of Answer Key}
\end{document}